% 调用上级目录公共模板
\documentclass{../template/softeng}

\begin{document}

% 封面：文档名、项目名、版本
\doccover{软件需求规格说明书}
{校园图书借阅管理系统}
{V2026.07}

% 目录
\tableofcontents
\newpage

% 修订记录
\begin{revrecord}
V1.0 & 2026-07-10 & 张三1 & 李四 & 初稿完成，完成功能性需求 \\
\hline
V1.1 & 2026-07-11 & 张三1 & 李四 & 完成非功能性需求梳理 \\
\hline
\end{revrecord}

% 正文正文结构（软件工程国标标准章节）
\section{引言}
\subsection{编写目的}
本文档为在线OJ（Online Judge）系统软件需求规格说明书，用于完整定义系统全部功能性与非功能性需求，明确项目开发边界、业务规则、数据规范、运行约束。
本文档交付对象包含项目经理、系统设计人员、前后端开发工程师、测试人员，作为项目规划、架构设计、代码开发、功能测试、验收交付的核心依据。

\subsection{项目背景}
本系统为高校大二软件工程实训项目，开发单位为在校实训开发小组，服务对象为高校程序设计课程教师与在校学生。
系统独立部署运行，无第三方外部系统强制对接，仅依托Web服务、判题沙箱、数据库完成完整业务闭环，用于支撑线上编程练习、自动代码评测、学情统计与教学管理。

\subsection{定义}
\begin{itemize}
    \item OJ：Online Judge，在线程序自动判题系统；
    \item AC：Accepted，代码通过全部测试用例；
    \item 判题沙箱：隔离运行用户代码的安全执行环境，防止恶意代码破坏服务器；
    \item 测试用例：包含标准输入、预期输出、边界条件的程序评测数据；
    \item 提交队列：待后台执行判题任务的任务缓冲队列；
\end{itemize}

\subsection{参考资料}
\begin{enumerate}
    \item GB/T 9385-2008 计算机软件需求规格说明规范；
    \item 实训项目任务书（2026软件工程实训）；
    \item Web前后端分离开发规范文档；
    \item 在线判题系统安全沙箱开发技术手册。
\end{enumerate}

\section{任务概述}
\subsection{目标}
搭建一套Web端在线编程判题平台，区分学生、教师两类角色；实现题库分级管理、在线代码编写提交、沙箱自动编译判题、提交记录查询、个人刷题排行、题目讨论交流；教师端提供题库批量维护、用户监管、平台运行监控、公告发布、讨论区管理功能，满足高校编程课程线上练习、自动批改、学情统计教学需求。

\subsection{运行环境}
\begin{itemize}
    \item 服务操作系统：Windows Server / Linux CentOS；
    \item Web服务：Nginx / Tomcat；
    \item 数据库：MySQL 8.0；
    \item 判题依赖：多语言编译器、隔离沙箱程序；
    \item 客户端：主流Chrome、Edge浏览器。
\end{itemize}

\subsection{条件与限制}
\begin{enumerate}
    \item 判题服务必须做代码隔离，禁止用户代码直接访问服务器本地文件；
    \item 系统需限制单用户单次代码运行时间、内存上限，避免服务器资源耗尽；
    \item 仅支持教师、学生两种角色，不开放游客完整答题权限；
    \item 题库批量导入仅支持标准格式化文本文件，不支持自定义复杂格式；
    \item 平台部署硬件资源有限，并发判题任务数量存在上限。
\end{enumerate}

\section{数据描述}
\subsection{静态数据}
系统长期持久存储不变基础数据：用户账号信息、题目基础信息、难度标签、测试用例、公告内容、讨论帖子、教师权限配置、系统基础参数。

\subsection{动态数据}
系统运行过程中实时产生、变更的数据：用户登录会话、代码提交记录、判题运行日志、实时在线人数、判题队列任务、刷题统计数据、排行榜实时数据、批阅标记记录。

\subsection{数据库介绍}
采用MySQL 8.0关系型数据库存储全部业务数据，划分用户表、题目表、测试用例表、提交记录表、讨论区表、系统公告表、系统监控数据表，数据表通过外键建立关联约束。

\subsection{数据词典}
包含用户ID、用户名、角色标识、题目ID、难度等级、输入输出样例、时间限制、内存限制、提交状态码、判题耗时、通过率、在线人数、队列长度等全部字段含义、数据类型、取值范围说明。

\subsection{数据采集}
\begin{itemize}
    \item 基础题库数据：教师手动录入或批量文件导入；
    \item 用户数据：学生注册、教师后台录入；
    \item 运行统计数据：系统后台定时采集在线、判题负载、提交数据；
    \item 交互数据：学生提交代码、发帖、评论行为实时采集入库。
\end{itemize}

\section{功能需求}
\subsection{功能划分}
系统整体划分为五大功能模块：学生端业务模块、教师题库管理模块、用户与提交监管模块、系统监控运维模块、讨论区与公告模块。

\subsection{功能描述}
\reqitem{FR001}{学生可注册账号、登录、修改个人信息，浏览分级题库。}
\reqitem{FR002}{题目按入门、简单、中等、困难划分，展示题面、输入输出、样例、时空限制。}
\reqitem{FR003}{学生在线编写代码提交，后台沙箱自动编译运行，返回判题结果。}
\reqitem{FR004}{学生查看个人提交记录、错误详情、解题数量与个人排名。}
\reqitem{FR005}{学生可进入题目讨论区发布思路、提问、回复他人帖子。}
\reqitem{FR006}{教师可新增、编辑、删除题目，配置多组测试用例与难度标签，批量导入题库。}
\reqitem{FR007}{教师查看全体学生账号、提交记录，标记抄袭等异常提交并限制账号。}
\reqitem{FR008}{教师可视化查看在线人数、判题队列、编译器负载、服务器健康状态。}
\reqitem{FR009}{教师发布全站公告，审核、删除讨论区违规内容。}
\reqitem{FR010}{系统自动生成全局刷题统计报表，支撑教师教学效果评估。}

\section{性能需求}
\subsection{数据精确度}
判题结果严格匹配测试用例标准输出，无字符、空格、换行误差；统计通过率、解题数量数据精确至整数，排行榜实时数据无统计偏差。

\subsection{时间特性}
\begin{itemize}
    \item 网页页面加载响应时间 ≤ 1.5s；
    \item 普通代码提交判题反馈时间 ≤ 3s；
    \item 监控数据刷新间隔 5s；
    \item 批量导入题库处理单条题目耗时 ≤ 0.1s；
\end{itemize}

\subsection{适应性}
系统可适配Windows、Linux服务端环境；支持新增编程语言编译器；可扩展班级分组统计功能；可调整判题并发上限、代码运行时空限制参数。

\section{运行需求}
\subsection{用户界面}
采用Web网页可视化界面，题库分页列表、在线代码编辑器、监控数据图表、表格化提交记录、分页讨论帖子；统一菜单布局、按钮样式，弹窗提示操作结果。

\subsection{硬件接口}
无需专用硬件外设接口，仅依赖服务器CPU、内存、磁盘资源完成代码编译与数据存储。

\subsection{软件接口}
数据库接口：MySQL标准JDBC/ODBC连接接口；
判题接口：本地进程调用各语言编译器；
前端接口：前后端RESTful API数据交互接口。

\subsection{故障处理}
\begin{itemize}
    \item 数据库断开：页面提示服务器异常，缓存用户未提交代码；
    \item 判题服务器宕机：任务存入队列，恢复后自动重试判题；
    \item 超长恶意代码：直接拦截，返回运行超限错误；
    \item 文件上传失败：弹窗提示并保留本地代码。
\end{itemize}

\section{其它需求}
\begin{enumerate}
    \item 安全性：区分角色权限隔离，沙箱隔离用户代码，防止SQL注入、恶意脚本攻击；
    \item 可维护性：后台可视化管理页面，日志完整记录所有提交、操作行为；
    \item 可移植性：模块化分层开发，可快速迁移至不同服务器环境；
    \item 易用性：操作流程简洁，新手学生可快速完成代码提交操作；
    \item 可扩展性：预留接口支持新增题型、新增编程语言、班级分组管理。
\end{enumerate}

% 插入图片示例（OJ系统用例图）
\begin{figure}[htbp]
    \centering
    \includegraphics[width=0.7\textwidth]{usecase.png}
    \caption{在线OJ系统整体用例图}
\end{figure}

% 代码块示例（简易判题结果状态枚举）
\begin{lstlisting}[language=Python]
# 判题结果状态枚举
class JudgeStatus:
    ACCEPT = "AC"          # 通过
    WRONG_ANSWER = "WA"    # 答案错误
    COMPILE_ERROR = "CE"   # 编译错误
    TIME_LIMIT = "TLE"     # 超时
    MEMORY_LIMIT = "MLE"   # 内存超限
    RUNTIME_ERROR = "RE"   # 运行时错误
\end{lstlisting}

\end{document}